\section{Finite logarithmic flags and angular paths}\label{sec:flags}

An unbounded sequence of equation arguments can involve several widely
separated scales. This section records those scales and their exact
logarithmic relations in a finite flag, then constructs complex paths
through the sequence anchors. Along these paths, the largest scale
approaches a vertical complex direction before the smaller scales do.
We construct nested contours on which their real parts remain positive
and quantitatively separated. These are the domains for the separation
argument in Section~\ref{sec:separation}. The final two lemmas control
anchor-dependent growth and simultaneous avoidance of finitely many
lattices of poles.


\subsection{Sequence conventions and exact logarithmic relations}

All constructions in this section concern a selected sequence of real
points. Passing to a subsequence is allowed. For positive quantities,
$X\gg Y$ means $X/Y\to+\infty$. A comparison involving finitely many
quantities, and comparisons of a finite list at every fixed iterated
exponential height, can be decided after a diagonal subsequence. None of
the resulting numbers of scales or heights is asserted to be uniform
over the original sequences.

The \emph{ordinary variables} have a finite limit $t_c$ in the interior of
an analytic argument neighborhood. Bounded ratios that a formula needs
are included among these arguments. Neighborhoods can be shrunk for
each finite request. The endpoint inequalities defining the actual
problem can nevertheless approach their boundary. An ordinary
coefficient equal to zero below means the zero analytic germ at $t_c$,
and not merely a function vanishing on the selected sequence.

\begin{definition}[Logarithmic flag]\label{def:flags-flag}
A logarithmic flag consists of positive scales
\[
 Z_1\gg Z_2\gg\cdots\gg Z_m\gg1,
 \qquad f_i=\log Z_i,\qquad h_i=e^{-Z_i},
\]
each declared either free or dependent. The tuple of free scales is
denoted by $S$. A dependent scale has the exact relation
\begin{equation}\label{eq:flags-triangular}
 f_i=\sum_{j>i}a_{ij}Z_j+b_i,
 \qquad
 \sigma(i)=\min\{j:a_{ij}\ne0\},\qquad a_{i,\sigma(i)}>0.
\end{equation}
The coefficients $a_{ij}$ are fixed real numbers, and the backgrounds
$b_i$ are bounded ordinary coordinates, or their current values on an
ordinary chart. The free scales are independent arguments; dependent
scales are always evaluated by Equation~\eqref{eq:flags-triangular}.

The saturation condition below will ensure that bounded changes of the
backgrounds preserve the ordered hierarchy. Two scales belong to the same
\emph{height block} if each is eventually
bounded by a finite iterated exponential of the other. A flag is
\emph{saturated} if, in every block containing a dependent node, all
dependent nodes precede all free nodes and
\begin{equation}\label{eq:flags-saturated}
 \log Z_{\text{first free}}=o\bigl(\log Z_{\text{last dependent}}\bigr).
\end{equation}
An all-free block has no additional condition.
\end{definition}

The relation defining height blocks is an equivalence relation: a
composition of finitely many iterated exponentials is another finite
iterate, after increasing the iterate to absorb fixed constants. Its
classes are consecutive in the ordering of the scales. A dependent
node and its leading child $\sigma(i)$ are in the same block, because
\[
 f_i=a_{i,\sigma(i)}Z_{\sigma(i)}(1+o(1)).
\]
In particular a dependency chain ends at a free node in the same block.

We use lifted logarithms and powers in complex continuations. Thus
$f_i$ continues as the expression in
Equation~\eqref{eq:flags-triangular}, even if $Z_i=e^{f_i}$ winds in
the plane. An exponential monomial is
\[
 h^\nu=\exp\Bigl(-\sum_{i=1}^m\nu_i Z_i\Bigr).
\]
Its order is the sign of the first nonzero entry of $\nu$, with the
largest scale first. This convention does not refer to the numerical
sign of the monomial. Known signs and ordinary units can be kept as
separate factors. Later expansions will be trees indexed by successive
prefixes of $\nu$; no simultaneous convergence of all these monomials
is implicit in this notation.

Inserting logarithms expresses fixed real powers of large inputs as
exponential monomials times bounded nonvanishing ordinary analytic factors. A later insertion of $-\log|t|$, when an ordinary
coordinate $t$ tends to zero, must also preserve the large coordinates
already free at that stage. The next lemma supplies these operations.

\begin{lemma}[Finite saturation and preservation of free coordinates]
\label{lem:flags-saturation}
A finite collection of large real inputs admits a flag after a fixed
linear change of coordinates and the addition of bounded ordinary
coordinates. One can insert the logarithm of any free scale, retaining
the old scale values, by making that scale dependent and adding at most
one positive free residual or one bounded ordinary residual. The free
residual is $O(\log S)$ if the scale being converted is $S$.

A flag admits a finite saturation by such insertions. After saturation,
promoting one ordinary coordinate $t\to0$ of known sign, by inserting
$-\log|t|$, admits a finite further saturation in which every old free
coordinate stays free. The number of ordinary coordinates does not
increase in this promotion.

\end{lemma}

\begin{proof}
For the initial linear reduction, choose a largest absolute value
among a basis of the finite span of the inputs. After a subsequence,
the ratios of the other elements to it have finite limits. Subtract
these limiting multiples. In the remaining span repeat the operation,
changing the sign of an unbounded residual when necessary. At each
step the dimension of the remaining span drops. A bounded residual is
retained as an ordinary coordinate, even if its selected values vanish; every unbounded residual chosen as a new scale
is strictly smaller than its predecessor. This gives a finite ordered
scale basis with constant coefficients.

Apply the same elimination to $\log S$ against the existing scale
basis, beginning with the largest scale. A scale much larger than
$\log S$ has zero coefficient. After subtracting its limiting multiple
at each comparable scale, the remainder is either bounded or, after
changing sign, a new positive scale separated from every old scale.
All the subtracted terms and the remainder are $O(\log S)$. Hence only
one new residual is needed. Since $\log S=o(S)$, its expression uses
only scales below $S$; it is therefore a relation of the form
Equation~\eqref{eq:flags-triangular}. Its first nonzero coefficient is
positive, since $\log S\to+\infty$. The old scale values have not been
changed.

Process the blocks containing dependent nodes from largest to
smallest. On entering a block fix its then smallest scale $T$.
Convert every free scale $R$ in that block for which
$\log R\ne o(\log T)$, and repeat this test only on any new residual
still in the block. Every continuing chain satisfies
\begin{equation}\label{eq:flags-residual-chain}
 R_{q+1}\le C_q\log R_q.
\end{equation}
The initial node of a chain is bounded by $\exp^{N}(T)$ for a fixed
$N$, by the definition of the block. Finitely many applications of
Equation~\eqref{eq:flags-residual-chain} give a residual $O(T)$ and
then one $O(\log T)$. The logarithm of the latter is
$o(\log T)$, so it is no longer converted. Fixed constants can be
absorbed by an extra logarithmic step. A residual that has entered a
lower block has already left the chain under consideration. There are
finitely many initial chains and each step has at most one child;
thus the procedure is finite. Residuals entering lower blocks are
free until that block is processed.

Every converted node in the block has, after the comparison
subsequence, $\log R\ge c\log T$. Every free node left at the end
has $\log F=o(\log T)$. The original dependent nodes also have
logarithms at least $\log T$. Consequently all dependent nodes
precede all remaining free nodes, and
Equation~\eqref{eq:flags-saturated} holds. All-free blocks need no
processing.

For a later promotion first insert $-\log|t|$ by the same linear
elimination. In each old block with dependent nodes keep the old cut
\[
 L=\log Z_{\text{old last dependent}}.
\]
Only the new residual chain is tested for conversion, with criterion
$\log R\ne o(L)$. For every old free node $F$,
$\log F=o(L)$. For each newly converted node $R$, the comparison
subsequence gives $\log R\ge cL$; therefore
\begin{equation}\label{eq:flags-old-free}
 \frac{\log F}{\log R}\longrightarrow0,
 \qquad \frac{R}{F}\longrightarrow+\infty.
\end{equation}
In particular a new residual below an old free node is never converted
in this block. Thus no old free coordinate is consumed, and the new
last-dependent cut still dominates the logarithms of all old free
nodes. The residual-chain argument proves termination here as well.
If the chain enters an old all-free block, leaving its residual free
introduces no dependent node into that block. Inserting new scales
cannot join two distinct old blocks: transitivity of finite-height
comparability would already have put their old members in one block.
The coordinate $t$ has been consumed; at most the final step of its
single chain produces a bounded residual. This proves the assertion
about the number of ordinary coordinates. The inverse expression for
$t$ has the form
\[
 t=\sgn(t)\exp\Bigl(-\sum_i\nu_iZ_i-B\Bigr),
\]
with $B$ bounded and with a positive first nonzero large coefficient
when such a coefficient occurs.

\end{proof}

The flag is now fixed. The next assertion concerns a different operation:
we hold its free inputs fixed and vary its bounded backgrounds, always
reevaluating dependent scales by their exact relations. This stability
will be needed when ordinary variables are changed implicitly.

\begin{lemma}[Stability under bounded backgrounds]
\label{lem:flags-stability}
On a saturated output flag of Lemma~\ref{lem:flags-saturation}, bounded
variations of the backgrounds at fixed free coordinates preserve the ordered scale hierarchy,
uniformly on each fixed bounded admissible real background set. More
precisely, let
\[
 H=\sum_{i=1}^m c_iZ_i+B,
\]
where the $c_i$ are fixed real numbers, their vector is nonzero,
and $B$ is uniformly bounded on that background set. If $k$ is the
first index with $c_k\ne0$, then
\[
 \frac{H}{c_kZ_k}\longrightarrow1
\]
uniformly under those variations; the denominator uses the current,
reevaluated value of $Z_k$. Thus the leading index and sign are
preserved. Dependent logarithms may enter such a comparison after
substitution of their exact triangular formulas. No sign assertion
is made when the resulting large part vanishes, or for an expression
containing an uninserted free logarithm.
\end{lemma}

\begin{proof}
Process the saturated output flag's height blocks from smallest to
largest. Each evaluated lower block is bounded above, uniformly on the fixed bounded real background set,
by a fixed finite exponential iterate of its largest free member.
Indeed it contains only finitely many triangular relations, whose
coefficients are fixed and whose backgrounds are bounded. Every free
member of a higher block dominates every such fixed iterate. Thus
lower-block values remain negligible compared with every free member
of the block being processed.

The free members of the current block are unchanged and ordered. If
the block has dependent members, let $d$ be its last one and put
$k=\sigma(d)$. The leading child lies in the same block and, by
saturation, is free. Its value is therefore unchanged. All later
scales have already been ordered, so the triangular relation gives
\[
 f_d=a_{dk}Z_k(1+o(1))
\]
uniformly on the background set. The logarithm of the first free
member is $o(f_d)$ on the original sequence. It remains so uniformly,
since that logarithm and $Z_k$ are fixed. This proves the
last-dependent/first-free boundary.

Now proceed upwards through the dependent members. The difference
of two consecutive dependent log relations is a fixed affine
combination of later scales and a bounded remainder. Its large part
is nonzero, since otherwise their logarithmic gap on the original
sequence would be bounded. Its first nonzero coefficient is positive
on that sequence. The later scales are already ordered uniformly by
the induction, so this coefficient still dominates. This proves all
remaining gaps and hence the full hierarchy uniformly.

Finally divide $H$ by $c_kZ_k$, evaluated at the current backgrounds.
Every later-scale ratio tends uniformly to zero, as does the bounded
remainder divided by $Z_k$. This proves the precise assertion about
affine comparisons. A combination involving dependent logarithms
first reduces to this form by exact substitution. If all large
coefficients cancel, their cancellation is an identity of fixed
coefficient vectors, leaving the current bounded remainder; it
supplies no additional sign conclusion.
\end{proof}

\subsection{Paths through far anchors}

We keep a saturated flag fixed and vary its free coordinates along a
complex path, with parameter $s=u+iv$. The path must pass through the
chosen anchor without
depending on the ordinary backgrounds, and it must preserve scale
separation on every required fixed backward interval. Its angular purpose
comes from the identity
\[
 \Re Z_i=e^{\Re f_i}\cos(\Im f_i).
\]
As $|v|$ grows, the larger scales must approach the directions
$\Im f_i=\pm\pi/2$ before the smaller ones. With
$\Lambda_i=f_i'$ on the real ray, the path estimates below give
$\Im f_i(u+iv)\sim v\Lambda_i(u)$: the level-$i$ direction is
therefore reached at distance about $\pi/(2\Lambda_i)$ from that
ray. We first construct the paths and estimate these logarithmic rates;
the next subsection chooses the precise contours and their margins.


Write $\operatorname{polylog}X$ for a bound
$C(1+\log X)^N$, where $C,N$ may change between occurrences and may
depend on a fixed derivative order. In every use the number of
derivatives and the other requests are finite. Thresholds, constants,
and ordinary neighborhoods may depend on this finite request.

Fix a far anchor, write its smallest logarithm as $f_m^0$, and put
$u_0=\log f_m^0$. We construct paths $s=u+iv$ through the free
coordinates, defined over $u\ge u_0-B$ for any required fixed
backstep $B$. All tuning constants are held fixed when testing a path.
The smallest node is necessarily free. Its path is
\[
 f_m(s)=e^s.
\]
Now take $i=m-1,\ldots,1$. Dependent nodes are evaluated from
Equation~\eqref{eq:flags-triangular}. At a free node $i<m$ set
$n=i+1$ and use the following rules.

If the successor $n$ is free, put $f_i=f_n+g_i$. A positive
logarithmic gap can usually be continued by
\begin{equation}\label{eq:flags-uncapped}
 g_i=x_i,\qquad x_i(s)=x_i^0 e^{s-u_0},
\end{equation}
with $x_i^0$ equal to the anchor gap. Within a block containing
dependent nodes, some free gaps must also stay small relative to the
leading child of the last dependent node. If that node is $d$ and
$k=\sigma(d)>i$, cap the gap by using
\begin{equation}\label{eq:flags-cap}
 g_i=\frac{x_i y_i}{x_i+y_i},\qquad
 y_i=\eta_i Z_k,
 \qquad
 x_i^0=\frac{g_i^0y_i^0}{y_i^0-g_i^0}.
\end{equation}
Here $k$ is free. Saturation gives $g_i^0=o(Z_k^0)$, so one can take
$\eta_i\to0$ so slowly that $g_i^0=o(y_i^0)$ and, for every
backstep needed in the finite request,
\begin{equation}\label{eq:flags-eta}
 \eta_i\ge Z_k(u_0-B)^{-1/4}.
\end{equation}
Then $x_i^0\sim g_i^0\to\infty$. The number $\eta_i$ is fixed
on the resulting path, though it tends to zero as anchors recede.

If the successor $n$ is dependent, saturation puts $i$ and $n$ in
different blocks. The new free path must dominate $f_n$ while remaining
independent of its backgrounds. For this purpose, follow leading children
from $n$ to its terminal free node
$k$, using $\ell$ edges, and put
\begin{equation}\label{eq:flags-boundary}
 T=\exp^{\ell-1}(Z_k^2),\qquad
 f_i=T^2+x_i,\qquad x_i=x_i^0e^{s-u_0}.
\end{equation}
The original distinct heights give $f_i^0\gg (T^0)^2$, so
$x_i^0=f_i^0-(T^0)^2>0$ and tends to infinity. Crucially, this rule
uses only the free node $k$, not an ordinary input or a dependent
value whose background could later change.

The real ordinary inputs can be constant, or can vary with bounded
derivatives of all the fixed orders used. Independent ordinary tests
are chosen real-symmetric and holomorphic, with positive-order
derivatives tending to zero, and return to $t_c$ as $u\to\infty$
on each individual path. For example, an anchor value $t^0$ in a
smaller convex box can be tested by
\[
 t(s)=t_c+(t^0-t_c)e^{-\delta(s-u_0)},
\]
with $\delta>0$ as small as needed and with slack for the fixed
backstep. This also shows why smaller ordinary boxes requested later
can be reached on the same ray. It is not necessary that all finite
requests hold at the anchor with one common threshold.

For complex estimates, a \emph{controlled background determination}
means a local continuation with uniform Taylor control on the angular
range in question. In particular it satisfies
\[
 b(u+iv)=b(u)+ivb'(u)+O(v^2),\qquad
 \partial_v b(u+iv)=ib'(u)+O(|v|),
\]
and the corresponding Taylor estimates through every fixed order
requested at that step. The real jets are bounded. Several piecewise
holomorphic determinations can have the same real jets; negligible
errors arising from their smooth interpolation are retained as defects
in the packet estimates below. Real boundedness alone is not used to
infer these complex bounds.

\begin{lemma}[Real and angular path estimates]\label{lem:flags-path}
The paths just constructed pass through the given anchors and
preserve the ordered scales under bounded controlled ordinary tests.
Put $\Lambda_i=f_i'$ on the real ray. For each $G=f_i-f_j$, $i<j$,
and for each fixed positive integer $r$, sufficiently far out,
\begin{equation}\label{eq:flags-real}
 \begin{gathered}
 G\longrightarrow+\infty,\qquad
 cG\le G'\le G\operatorname{polylog}G,\\
 |G^{(r)}|\le G\operatorname{polylog}G,\qquad
 0<\frac{(f_i-f_{i+1})'}{f_i-f_{i+1}}
       \le C\Lambda_{i+1}.
 \end{gathered}
\end{equation}
The first three bounds also hold with $G$ replaced by $f_i$.
In particular
\begin{equation}\label{eq:flags-rates}
 \Lambda_1>\cdots>\Lambda_m\longrightarrow+\infty,
 \qquad cf_i\le\Lambda_i\le f_i\operatorname{polylog}f_i,
 \qquad \frac{f_i''}{\Lambda_i^2}=o(1).
\end{equation}
Fixed backsteps can reduce $f_i$ by any prescribed fixed factor.

With holomorphic background tests, the formulas for indices $j\ge p$
have controlled holomorphic continuations on $|v|\le C/\Lambda_p(u)$
and satisfy
\begin{equation}\label{eq:flags-angular}
 \begin{aligned}
 \Re f_j(u+iv)&=f_j(u)+o(1),\\
 \Im f_j(u+iv)&=(1+o(1))v\Lambda_j(u),\\
 \Im(f_j-f_l)(u+iv)&=(1+o(1))v(\Lambda_j-\Lambda_l)(u),
                       && l>j,\\
 \partial_v\Im f_j(u+iv)&=(1+o(1))\Lambda_j(u).
 \end{aligned}
\end{equation}
These estimates also hold for controlled local determinations of the
backgrounds that agree with the real path to their requisite finite
Taylor orders; their antiholomorphic defects, if present, are accounted
for separately. For a retained free input $Z_l$, every fixed path
derivative is bounded by $\exp(C_r f_l(u))$ on its applicable angular
range.
\end{lemma}

\begin{proof}
We give the induction estimates, including the cases in which a
smaller scale can have a relatively large logarithmic derivative.
The initial function $f_m=e^u$ has all the required estimates.
If a positive function $F$ satisfies the first three real bounds,
the formula for a derivative of $e^F$ is $e^F$ times a polynomial in
$F',\ldots,F^{(r)}$. Consequently
\begin{equation}\label{eq:flags-exp-derivative}
 |(e^F)^{(r)}|\le e^F\operatorname{polylog}(e^F).
\end{equation}
Finite sums with a separated positive leading term obey the same
upper bounds. Their first derivative has that leading positive term:
the derivative ordering already known at later nodes and the
exponential modulus gaps make the lower terms negligible. This
handles dependent logs and differences of dependent logs. Fixed
bounded derivatives of the backgrounds are lower-order terms.

An uncapped gap has $x_i^{(r)}=x_i$ and poses no difficulty. For a
cap write $x=x_i$, $y=y_i$, $r_0=x/y$. Direct differentiation gives
\begin{equation}\label{eq:flags-cap-derivative}
 \frac{g'}g=\frac{1+r_0\Lambda_k}{1+r_0}.
\end{equation}
This is positive, bounded below by a fixed positive constant, and
bounded above by $1+\Lambda_k\le C\Lambda_n$. In addition,
Equation~\eqref{eq:flags-eta} gives
$\Lambda_k\le\operatorname{polylog}y$. If $x\le y$, then
$g\asymp x$, and the elementary bound
\[
 \frac{x}{y}(1+\log y)^N\le C_N(1+\log x)^N
 \quad(x\le y,\ x\hbox{ sufficiently large})
\]
absorbs each occurrence of a $y$ derivative. To verify the bound,
write $y=xe^a$ and use
$e^{-a}(1+\log x+a)^N\le C_N(1+\log x)^N$.
If $x\ge y$, then $g\asymp y$ and the bound is immediate.
Repeated differentiation of $xy/(x+y)$ gives the same alternatives:
terms containing a derivative of $y$, after division by $g$, have a
factor $O(x/y)$ when $x\le y$, and otherwise have only fixed powers
of logarithms of $y$. This proves the fixed-order bounds in $g$.

For the boundary rule, repeated use of the positive leading terms
on the chain from $n$ to $k$ gives, sufficiently far out,
\begin{equation}\label{eq:flags-height-offset}
 \exp^{\ell-1}(Z_k^{1/2})\le f_n
       \le \exp^{\ell-1}(Z_k^2)=T.
\end{equation}
Every fixed logarithmic derivative of $T$ is bounded by a power of
$1+\log T$. Its first logarithmic derivative is $o(f_n)$: for
$\ell=1$ this follows from
$T'/T=2\Lambda_k=o(Z_k^{1/2})$; for larger $\ell$, the numerator
is a product of lower iterates, whereas the left side of
Equation~\eqref{eq:flags-height-offset} has one more exponential
height. Thus $T^2+x_i$ and its first derivative dominate $f_n$ and
$f_n'$. The gap $T^2+x_i-f_n$ has the first three real estimates,
and
\[
 \frac{(T^2+x_i-f_n)'}{T^2+x_i-f_n}
 \le C\left(1+\frac{T'}T\right)\le C\Lambda_n.
\]
The same computations give its fixed-order bounds.

There remains the dependent-to-free boundary in a saturated block.
Let $d$ be its last dependent node and $k=\sigma(d)$. The only
free gaps above $k$ are the caps in
Equation~\eqref{eq:flags-cap}, so their finite sum is $o(Z_k)$ at
the anchors, with first derivative $o(Z_k\Lambda_k)$. For example,
Equation~\eqref{eq:flags-cap-derivative} and $g\le\eta_iZ_k$
bound the derivative by $C\eta_iZ_k(1+\Lambda_k)$.
On each individual forward ray, $x_i/Z_k\to0$, because $x_i$ grows
as $e^u$ and $f_k\ge e^u$ eventually. Hence the same little-oh
bounds hold towards infinity on that ray, with its tuning constants
fixed. The remaining free contribution is $f_k=o(Z_k)$, whose
first derivative is $o(Z_k\Lambda_k)$. Therefore
\[
 f_{d+1}=o(Z_k),\qquad
 f_{d+1}'=o(Z_k\Lambda_k),
\]
whereas $f_d\sim a_{dk}Z_k$ and
$f_d'\sim a_{dk}Z_k\Lambda_k$.
This proves the real estimates across this boundary, including
\[
 \frac{(f_d-f_{d+1})'}{f_d-f_{d+1}}
       \asymp\Lambda_k\le\Lambda_{d+1}.
\]
Higher derivatives of the caps are bounded by
$Z_k\operatorname{polylog}Z_k$, which suffices here. Between two
dependent nodes the leading scale of their difference has index at
least that of the leading scale in the earlier relation. Its
logarithmic derivative is at most $C\Lambda_{i+1}$: if its index
is $i+1$ this is immediate, and otherwise use the already proved
ordering of later derivatives. This completes the adjacent-gap
induction. A nonadjacent gap is a finite sum of adjacent positive
gaps, proving its first three estimates. The lower constant can,
for example, be taken to be $c=1/2$ after going sufficiently far
for any chosen fixed backstep: uncapped terms and caps have
logarithmic derivative at least one, and dependent leading terms
and the $T$ terms have logarithmic derivatives tending to infinity.
The discarded terms are relatively negligible. Thus a backstep
$B\ge2\log R$ reduces every $f_i$ by at least a prescribed factor
$R$, once that fixed backstep is included in the construction.

For completeness, the continuation uses the explicit formulas, not
an arbitrary holomorphic function having prescribed real bounds.
In the downward induction, a dependent $f_j$ is a finite affine sum
of $Z_k$ with $k>j$. The already controlled real parts of $f_k$
make $|Z_k(u+iv)|\asymp Z_k(u)$, so differentiating this sum
retains the bounds of Equation~\eqref{eq:flags-exp-derivative}.
The same argument applies directly to a dependent gap, using its
first surviving coefficient. In a boundary term $T$, every proper
exponent in its finite exponential tower and its fixed derivatives
are bounded by a power of $1+\log T$. Since
$|v|\le C/\Lambda_i\le C'/T^2$, their variations tend to zero.
Induction through that finite tower therefore preserves its moduli
and the derivative estimates. An uncapped additive term has its
explicit exponential continuation.
In a cap, if $|x|$ and $|y|$ are separated, $x+y$ stays away from
zero. If they are comparable, then $y=O(f_i)$ at the controlling
node, and
$|\arg y|=O(\Lambda_k/\Lambda_i)=o(1)$; indeed in this case
$\Lambda_i\ge cg\asymp y$, whereas
$\Lambda_k\le\operatorname{polylog}y$.
The denominator therefore again stays away from zero. These
observations prove by the same downward induction that the real
derivative bounds used above remain valid, with fixed enlarged
constants, along the short vertical segment.

Taylor's formula along that segment now gives
\[
 |\Re f_j(u+iv)-f_j(u)|
 \le C v^2 f_j(u)\operatorname{polylog}f_j(u)=o(1),
\]
since $\Lambda_p\ge cf_p\ge cf_j$. For a gap $G=f_j-f_l$
it gives an error bounded by
$Cv^2G\operatorname{polylog}G$. Dividing by
$|v|G'\ge c|v|G$ gives
$O(\operatorname{polylog}G/\Lambda_p)=o(1)$.
Applying this estimate directly to the gap, rather than subtracting
two relative estimates for individual logs, proves the third line
of Equation~\eqref{eq:flags-angular} even for nearly equal speeds.
The individual imaginary and derivative assertions follow in the
same way. Controlled backgrounds have the same Taylor bounds and
therefore give the stated version for local determinations. Finally,
the derivative polynomial for $Z_l=e^{f_l}$ is bounded by
$e^{f_l}\operatorname{polylog}(e^{f_l})\le e^{C_rf_l}$.
\end{proof}

\begin{remark}\label{rem:flags-path-heights}
Different original height blocks need not stay different along an
individual forward ray. The conclusions that persist are the ordered
scale ratios, the exact logarithmic relations, and
Equations~\eqref{eq:flags-real}--\eqref{eq:flags-angular}. These
are the conclusions used below. The paths of free inputs were defined
without ordinary backgrounds, which permits independent ordinary
tests and later regular implicit substitutions.
All comparisons defining a fixed finite real regime have slack
sufficiently far along the selected sequence. Continuity of the
finite triangular formulas therefore permits sufficiently small
absolute perturbations of the free inputs while retaining that
regime and the limiting ordinary data. Their permitted sizes may
depend on the anchor; no fixed neighborhood is intended.
\end{remark}

\subsection{Nested contours and background variation}

At level $p$ we use two boundaries before $Z_p$ reaches a vertical
direction. A good contour leaves $\Re Z_p$ of order
$V_p=e^{f_p}/\Lambda_p$, large enough to dominate the later real
parts. The outer side is closer to vertical; its defining barrier
retains both that domination and an explicit polynomial angular margin.
These margins also determine how accurately two background determinations
must agree.

\begin{lemma}[Good contours and buffered sides]\label{lem:flags-contours}
For a sufficiently small fixed $\eps>0$, the upper and lower
$p$-good contours are defined by
\begin{equation}\label{eq:flags-good}
 |\Im f_p|\sim\frac\pi2,\qquad
 \cos(\Im f_p)=\frac{\eps}{\Lambda_p(u)}.
\end{equation}
They have collars of fixed small multiples of their cosine margin
and are nested in increasing order of $p$. Put
\[
 V_p(u)=\frac{e^{f_p(u)}}{\Lambda_p(u)}.
\]
On a $p$-good contour, for $j>p$ and $G=f_p-f_j$,
\begin{equation}\label{eq:flags-good-real}
 \Re Z_p\sim\eps V_p,
 \qquad
 \Re Z_j\asymp
 \frac{e^{f_j(u)}(\eps+G'(u))}{\Lambda_p(u)}
       =o(V_p).
\end{equation}

For large fixed $A$ and fixed $K\ge2$, an outer side can be placed
inside the signed vertical direction by the barrier
\begin{equation}\label{eq:flags-side}
 L_p=Z_p-A\sum_{j>p}Z_j-\frac{Z_p}{f_p^K},
 \qquad \Re L_p=0.
\end{equation}
We use its inward component, where $\Re L_p>0$. Throughout this
component, with harmless fixed adjustments near its boundary,
\begin{equation}\label{eq:flags-buffer}
 \Re Z_p\ge A\sum_{j>p}\Re Z_j
               +c\frac{e^{f_p(u)}}{f_p(u)^{K+1}},
 \qquad \Re Z_p\gg f_p(u)^N
 \quad\hbox{for every fixed }N.
\end{equation}
The side and good contours can be chosen with extra room before any
finite subsequent narrowing. The same conclusions hold if the
quantities in the barrier are replaced by holomorphic proxies with
absolute $o(1)$ errors.

Uniformly between the $p$-good contour and its outer side, for
$p<n<l$,
\begin{equation}\label{eq:flags-later-real}
 \Re Z_l=o(\Re Z_n).
\end{equation}
Every retained index $j\ge p$ has a polynomial buffer of its own:
for some fixed exponent $C$,
\begin{equation}\label{eq:flags-retained-buffer}
 \Re Z_j\ge\frac{e^{f_j(u)}}{(1+f_j(u))^C}.
\end{equation}
In particular $\Re Z_j/f_l(u)\to\infty$ for $l\ge j$, so
fixed derivative losses $\exp(C'f_l)$ can be absorbed by any fixed
positive exponential decay in $\Re Z_j$.
\end{lemma}

\begin{proof}
By Equation~\eqref{eq:flags-angular}, $\Im f_p$, the lifted phase
of $Z_p$, is strictly increasing with $v$ near the positive
vertical position of $Z_p$,
which is at $v=(1+o(1))\pi/(2\Lambda_p)$. This gives the upper
contour, and real symmetry gives the lower one. If $j>p$, the
relative gap formula, with
$\eta_p=\cos(\Im f_p)\ge0$, gives
\begin{equation}\label{eq:flags-cosine}
 \Re Z_j\asymp e^{f_j(u)}
   \left(\eta_p+
               \frac{\Lambda_p-\Lambda_j}{\Lambda_p}\right)
\end{equation}
near the vertical direction. This also holds with the evident
constant-factor interpretation when the ratio of speeds is bounded
away from one. On a good contour it proves
Equation~\eqref{eq:flags-good-real}, since
$e^{-G}(\eps+G')\to0$ by the polylogarithmic upper bound for
$G'$. Moreover $G'\to\infty$, so the later phase is farther from
vertical than any fixed-width $p$ collar. This proves nesting even
when the two speeds have ratio tending to one. Farther inside, all
relevant cosines are bounded below, and the same ordering is simpler.

To compute the buffer, write $f_p=a+i\theta$, with
$a=f_p(u)+o(1)$ and $|\theta|\le\pi/2$. Uniformly near the upper
vertical direction,
\begin{equation}\label{eq:flags-buffer-expansion}
 \Re\frac{Z_p}{f_p^K}
 =e^a a^{-K}\left(
   \cos\theta+\frac{K\theta}{a}\sin\theta+O(a^{-2})\right).
\end{equation}
The second term is positive and of size $1/a$ near vertical.
The lower direction gives the same sign, since
$\theta\sin\theta>0$. Away from vertical, the first term gives
an even larger lower bound. Thus the real part in
Equation~\eqref{eq:flags-buffer-expansion} is at least
$ce^a/a^{K+1}$ throughout the required wedge.
At the good contour this term and the later real phases are
$o(\Re Z_p)$, so $\Re L_p>0$ there. At the exact vertical
direction $\Re Z_p=0$ and both subtracted real terms are positive.
The inward component therefore meets a zero side before vertical,
and positivity of its barrier gives
Equation~\eqref{eq:flags-buffer}. More explicitly, division of the
side equation by $e^{f_p(u)}$, followed by
Equation~\eqref{eq:flags-cosine}, gives, with $F=f_p(u)$,
\[
 \eta_{p,\mathrm{side}}\asymp F^{-K-1}
   +A\sum_{j>p}e^{-(f_p-f_j)}
                    \frac{\Lambda_p-\Lambda_j}{\Lambda_p}.
\]
The terms proportional to $\eta_{p,\mathrm{side}}$ on the
right have coefficient tending to zero and were absorbed on the
left. Hence $\Lambda_p\eta_{p,\mathrm{side}}\to0$, by the
polylogarithmic gap estimates. The side therefore lies beyond
every fixed good contour and still before the vertical direction.
This argument uses no global
injectivity of $Z_p$ or of a proxy. An initial choice with smaller
$A$ and larger $K$ supplies extra room; increasing $A$ and slightly
decreasing $K$ afterwards moves the working side inward. Absolute
$o(1)$ corrections are negligible compared with its diverging
buffer. The identical sign argument constructs the corrected side
within that extra room.

For the relative statement let $H=f_n-f_l$. From
Equation~\eqref{eq:flags-cosine},
\[
 \frac{\Re Z_l}{\Re Z_n}
 \le C e^{-H}\left(
  1+\frac{\Lambda_n-\Lambda_l}
  {\Lambda_p\eta_p+\Lambda_p-\Lambda_n}\right).
\]
The denominator is at least
$(f_p-f_n)'\to\infty$, whereas
$\Lambda_n-\Lambda_l\le H\operatorname{polylog}H$.
The right side tends to zero uniformly, proving
Equation~\eqref{eq:flags-later-real}. This is why one fixed
increase of $A$ can pay any prescribed finite collection of later
phase losses at the current side.

Finally consider $j>p$. If
$\Lambda_j\le\Lambda_p/2$, Equation~\eqref{eq:flags-cosine}
gives a fixed positive cosine for $Z_j$. Otherwise
$\Lambda_p\le2\Lambda_j$, and
\[
 f_p\le C\Lambda_p\le C f_j\operatorname{polylog}f_j.
\]
The polynomial angular margin supplied by
Equation~\eqref{eq:flags-buffer} for $p$ is therefore also an
inverse power of $f_j$. Equation~\eqref{eq:flags-cosine} again
proves Equation~\eqref{eq:flags-retained-buffer}. The case $j=p$
is already Equation~\eqref{eq:flags-buffer}.
\end{proof}

We next compare the domains obtained from different determinations of
the same real backgrounds. At a level-$p$ contour, first-jet agreement
preserves the separation of later phases from the vertical direction.
Locating the perturbed determination's own thin side requires higher,
but still finite, agreement.

\begin{lemma}[Sensitivity to background jets]\label{lem:flags-background}
Let $q\ge2$ be a fixed integer. On
$|v|\le C/\Lambda_p(u)$, two controlled background
determinations differing by $O(|v|^q)$, and agreeing on the real
path, change any $f_l$, $l\ge p$, by
\begin{equation}\label{eq:flags-background}
 |\Delta f_l|\le C\Lambda_l(u)|v|^q.
\end{equation}
The estimate is uniform along the short segment between the two
determinations when that segment is taken with the same background
slack. In particular agreement through the real first jet gives
$O(\Lambda_l v^2)$. On the angular scale $|v|=O(1/\Lambda_p)$
this error is smaller than the phase gain $(\Lambda_p-\Lambda_l)/\Lambda_p$ for $l>p$.
Agreement through a sufficiently high fixed jet preserves any
prescribed polynomial angular buffer and good-contour collar.
\end{lemma}

\begin{proof}
A free log has zero background sensitivity. For a dependent node,
differentiate Equation~\eqref{eq:flags-triangular}. By downward
induction and the leading positive derivative estimate,
\[
 \sum_{k\ge\sigma(l)}Z_k(u)(1+\Lambda_k(u))
       \le C\Lambda_l(u).
\]
Exponential moduli remain comparable on the stated angular scale.
The derivative of $f_l$ with respect to each bounded background is
therefore at most $C\Lambda_l(u)$. Integrating this derivative
between the two determinations proves
Equation~\eqref{eq:flags-background}; it also proves inductively
that the intermediate determinations remain in the allowed range.
For $q=2$ the error is at most $C/\Lambda_p$, while
$(\Lambda_p-\Lambda_l)/\Lambda_p\gg1/\Lambda_p$ by
Equation~\eqref{eq:flags-real}. For an own-contour comparison,
$\Lambda_p|v|^q=O(\Lambda_p^{1-q})$. The two-sided polynomial
bounds for $\Lambda_p$ in $f_p$ imply that choosing a fixed $q$
large enough makes this smaller than each of the finitely many
specified polynomial margins. No infinite-order jet matching is
being asserted.
\end{proof}

When backgrounds have slightly different band determinations, the
first-jet comparison in Lemma~\ref{lem:flags-background} suffices
for later phase comparisons: its $O(1/\Lambda_p)$ angular error
is swallowed by the diverging gap derivative. Locating a
determination's own thin side or its own good collar by a common
reference uses the higher finite jet comparison instead. This
distinction will be retained in ordinary implicit substitutions.

\subsection{Exceptional growth and safe columns}

An uncapped summand can make $\log(2+\Lambda_p)$ too large to
absorb in $\Re Z_{p+1}$. The constant $D_p$ below records the
remaining anchor dependence. The second lemma selects vertical columns
that avoid finitely many real pole lattices; when $D_p$ is large, it
also preserves this avoidance throughout the collar windows needed in
separation. All these assertions concern a fixed finite request.

\begin{lemma}[The exceptional anchor constant]\label{lem:flags-exceptional}
For $p<m$ set $D_p=\log(2+x_p^0)$ if $f_p$ contains an
uncapped slow additive term in
Equation~\eqref{eq:flags-uncapped} or
Equation~\eqref{eq:flags-boundary}, and set $D_p=1$ otherwise.
Set $D_0=1$. On the buffered level-$n$ wedge, $n=p+1$,
\begin{equation}\label{eq:flags-D}
 \log(2+\Lambda_p)\le C D_p+o(\Re Z_n).
\end{equation}
If $D_p$ is large, then throughout the fixed-backstep range
\begin{equation}\label{eq:flags-D-speed}
 \Lambda_p(u)\ge c_B e^{D_p}e^{u-u_0}.
\end{equation}
\end{lemma}

\begin{proof}
If $p$ is dependent, $f_p=O(Z_n)$ and
$\Lambda_p\le Z_n\operatorname{polylog}Z_n$ by its triangular
relation. Thus $\log(2+\Lambda_p)=O(f_n+\log f_n)$, which is
$o(\Re Z_n)$ by the polynomial buffer. A cap has size at most
a later scale and its logarithmic derivative is bounded by a
polylogarithm of that scale, giving the same conclusion. In the
boundary construction, the height offset gives
\[
 \log T=o\left(\frac{Z_n}{f_n^N}\right)
 \quad\hbox{for every fixed }N;
\]
this follows directly from
Equation~\eqref{eq:flags-height-offset}, since one additional
exponential dominates all its fixed powers. It bounds the
$T$ contribution to $\log(2+\Lambda_p)$. The only remaining
contribution is
\[
 \log(2+x_p(u))\le D_p+O(1+|u-u_0|).
\]
The second term is $o(\Re Z_n)$ on the forward ray and on each
fixed backstep, because $f_n\ge f_m=e^u$ and the buffer grows
exponentially in $f_n$. This proves
Equation~\eqref{eq:flags-D}. An uncapped term contributes exactly
$x_p^0e^{u-u_0}$ to the positive derivative of $f_p$; its other
terms have positive derivatives. For large $D_p$,
$x_p^0\asymp e^{D_p}$, proving
Equation~\eqref{eq:flags-D-speed}.
\end{proof}

The following elementary consequence is recorded for use when raw
coefficients have poles at finitely many real lattices. It asserts
safety for each finite request, not for an infinite set of labels.

\begin{lemma}[Common safe columns]\label{lem:flags-safe-columns}
Let a fixed finite list of positive real sweeps $x_a(u)$ and real
lattices $c_a+h_a\ZZ$, $h_a>0$, satisfy, uniformly far out,
\begin{equation}\label{eq:flags-sweeps}
 \begin{gathered}
 x_a\to\infty,\qquad x_a'>0,\qquad
 \inf_{u\ge U}x_a'(u)\longrightarrow\infty\quad(U\to\infty),\\
 \left|\frac{x_a''}{(x_a')^2}\right|\le\frac{C_a}{x_a}.
 \end{gathered}
\end{equation}
For a sufficiently small fixed $\delta>0$, every fixed-length
real interval far out contains a common column $u_*$ with
\[
 \dist(x_a(u_*),c_a+h_a\ZZ)\ge2\delta
 \quad\hbox{for all }a.
\]
If their complex continuations obey
\[
 |\Im x_a(u_*+iv)|<1
 \quad\Longrightarrow\quad
 |\Re x_a(u_*+iv)-x_a(u_*)|=o(1),
\]
these columns stay a fixed positive distance from all the listed
real poles on their whole applicable angular ranges.

Suppose additionally that on a band $j=p+1$ the relevant sweeps
satisfy $x_a'=O(Z_j\Lambda_j)$ and that at the chosen column
\begin{equation}\label{eq:flags-D-active}
 D_p>\frac{Z_j(u_*)}{f_j(u_*)^{K+3}}.
\end{equation}
Then safety persists throughout every required fixed collar window
\[
 |f_p(u)-f_p(u_*)|\le C_1\log(2+D_p)
\]
in the fixed-backstep range.
\end{lemma}

\begin{proof}
For one sweep use $x=x_a$ as a coordinate and let
$w(x)=1/x_a'(u(x))$. Then
\[
 \frac{\dd}{\dd x}\log w(x)
       =-\frac{x_a''}{(x_a')^2}=O(1/x).
\]
On any fixed lattice period the ratio of the largest and smallest
values of $w$ tends to one. The time spent within distance
$2\delta$ of its lattice is consequently at most
$C_a'\delta/h_a$ times the period's total time. Sum over the
complete periods of a fixed $u$-interval. At most two incomplete
end periods remain; their total time tends to zero by the lower
bound on $x_a'$. Thus the bad proportion for this sweep is at most
$C_a'\delta/h_a+o(1)$. A union bound over the finite list makes
the common bad proportion less than one, by first choosing
$\delta$ sufficiently small. This does not require independent
or incommensurable sweeps. If the complex imaginary part has
absolute value at least one, it already separates the point from
the real lattice. Otherwise the assumed real shift is $o(1)$,
which preserves the original detuning with slack.

Under Equation~\eqref{eq:flags-D-active},
\[
 \log(Z_j\Lambda_j)=O(f_j+\log f_j)=o(D_p).
\]
Equation~\eqref{eq:flags-D-speed} makes the stated $f_p$ window
have real width $O(e^{-D_p}\log(2+D_p))$. On this window the
real derivative bounds ensure the same local estimates for
$Z_j\Lambda_j$; for example its logarithmic derivative has
size at most $\operatorname{polylog}Z_j=\exp(o(D_p))$ and
therefore changes its logarithm by $o(1)$. Each sweep changes by
at most
\[
 e^{-D_p+o(D_p)}\log(2+D_p)=o(1).
\]
The detuning again persists with slack. This proves the collar
assertion.
\end{proof}
